% =====================================================================
%  PIVOT — identifiability analysis (T1/T2/T3 of my_work/THEORY.md)
%
%  논문 3_method.tex 의 Problem Formulation 바로 뒤에 붙이면 됩니다.
%  전제 패키지: amsmath, amssymb, amsthm  (IEEEtran 에서 아래 선언 필요)
%
%    \newtheorem{lemma}{Lemma}
%    \newtheorem{proposition}{Proposition}
%    \newtheorem{corollary}{Corollary}
%    \newtheorem{remark}{Remark}
%
%  모든 결과는 my_work/study_theory.py 가 수치로 검증합니다 (상대오차 8e-16).
% =====================================================================

\subsection{Measurement Model and Notation}
\label{sec:model}

The object consists of $P$ rigid bodies (parts and hinges alike); body $i$ has
volume $V_i$, obtained from the reconstruction, and unknown density $\rho_i$.
The object is grasped rigidly, and $\mathcal{S}$ denotes the wrist force/torque
sensor frame. At articulation configuration $\theta\in\Theta\subset\mathbb{R}^J$
the centroid of body $i$ expressed in $\mathcal{S}$ is $c_i(\theta)$, which
follows from forward kinematics and geometry alone and is therefore independent
of $\rho$. Writing $\hat g\in\mathbb{S}^2$ for the gravity direction in
$\mathcal{S}$ and $g$ for its magnitude, the quasi-static wrench is
%
\begin{equation}
f = g\Big(\textstyle\sum_i V_i\rho_i\Big)\hat g,
\qquad
\tau = g\sum_i V_i\rho_i\,\big(c_i(\theta)\times\hat g\big).
\end{equation}
%
Collecting the geometry into
%
\begin{equation}
v \triangleq (V_1,\dots,V_P)^{\!\top}\!\in\mathbb{R}^{P},
\qquad
C(\theta) \triangleq \big[\,V_1c_1(\theta)\;\cdots\;V_Pc_P(\theta)\,\big]
\in\mathbb{R}^{3\times P},
\end{equation}
%
and using $c\times\hat g=-[\hat g]_\times c$, the wrench is linear in $\rho$:
%
\begin{equation}
y = A(\theta,\hat g)\,\rho + \varepsilon,
\qquad
A(\theta,\hat g) = g\begin{bmatrix}\hat g\,v^{\!\top}\\[2pt]
-[\hat g]_\times C(\theta)\end{bmatrix}\in\mathbb{R}^{6\times P}.
\label{eq:regressor}
\end{equation}
%
The sensor noise is block-isotropic,
$R=\mathrm{blkdiag}(\sigma_f^2 I_3,\;\sigma_\tau^2 I_3)$, which holds for the
per-axis specification of our wrist sensor. The robot reorients the wrist to
realize a set $\mathcal{G}=\{\hat g_1,\dots,\hat g_K\}$ of gravity directions
per configuration; we use an orthonormal triad ($K=3$).

\subsection{The Information Matrix in Closed Form}
\label{sec:lemma}

\begin{lemma}[Closed form]
\label{lem:closed}
For any finite $\mathcal{G}=\{\hat g_k\}_{k=1}^{K}$ of unit vectors, the Fisher
information about $\rho$ collected at configuration $\theta$ is
%
\begin{equation}
M(\theta;\mathcal{G})
= g^2\!\left[\frac{K}{\sigma_f^2}\,v v^{\!\top}
+\frac{1}{\sigma_\tau^2}\,C(\theta)^{\!\top}
\Big(K I_3-\!\sum_k \hat g_k\hat g_k^{\!\top}\Big) C(\theta)\right].
\label{eq:closed}
\end{equation}
\end{lemma}

\begin{IEEEproof}
From \eqref{eq:regressor},
$A^{\!\top}\!R^{-1}A
= g^2\big[\sigma_f^{-2}\,v(\hat g^{\!\top}\hat g)v^{\!\top}
+\sigma_\tau^{-2}\,C^{\!\top}[\hat g]_\times^{\!\top}[\hat g]_\times C\big]$.
Since $\|\hat g\|=1$ we have $\hat g^{\!\top}\hat g=1$, and
$[\hat g]_\times^{\!\top}[\hat g]_\times=-[\hat g]_\times^2
=I_3-\hat g\hat g^{\!\top}$. Summing over $k$ gives \eqref{eq:closed}.
\end{IEEEproof}

\begin{corollary}[Triad invariance]
\label{cor:invariance}
If $\mathcal{G}$ is an orthonormal basis of $\mathbb{R}^3$, then
$\sum_k\hat g_k\hat g_k^{\!\top}=I_3$ and
%
\begin{equation}
M(\theta)=g^2\!\left[\frac{3}{\sigma_f^2}v v^{\!\top}
+\frac{2}{\sigma_\tau^2}C(\theta)^{\!\top}C(\theta)\right],
\label{eq:triad}
\end{equation}
%
which is \emph{independent of the orientation of the triad}: rotating
$\mathcal{G}$ by any $Q\in SO(3)$ leaves $M(\theta)$ unchanged.
\end{corollary}

Corollary~\ref{cor:invariance} collapses the design space from
$\Theta\times SO(3)$ to $\Theta$: how the wrist is tilted carries no
information, and the articulation configuration is the only design variable.
Two further consequences follow directly from \eqref{eq:triad}.

\begin{remark}[The force channel is redundant with a scale]
\label{rem:force}
The force contribution is the rank-one term $\propto v v^{\!\top}$ for every
$\hat g$; it constrains only $v^{\!\top}\!\rho$, which is exactly the total mass
already given by a scale. All part-wise information therefore comes from the
torque channel, making $\sigma_\tau$ --- not $\sigma_f$ --- the binding sensor
specification.
\end{remark}

\begin{remark}[Why the loop must be closed]
\label{rem:closed-loop}
$M(\theta)$ in \eqref{eq:triad} does not depend on $\rho$: under exact angle
measurement the optimal configuration sequence could be computed offline, and
feedback would be unnecessary. Closing the loop is required \emph{only} because
$\theta$ is measured with error, which enters through the effective covariance
$R_{\mathrm{eff}}=R+J\Sigma_\theta J^{\!\top}$ with
$J=\partial(A\rho)/\partial\theta$, and hence depends on the current estimate.
\end{remark}

\subsection{Identifiability}
\label{sec:identifiability}

\begin{proposition}[Rank ceiling]
\label{prop:rank}
For any $K$ and any arrangement of gravity directions,
%
\begin{equation}
\operatorname{rank}M(\theta;\mathcal{G})\;\le\;1+\operatorname{rank}C(\theta)
\;\le\;1+\min(3,P).
\end{equation}
\end{proposition}

\begin{IEEEproof}
$\hat g_k\hat g_k^{\!\top}\preceq I_3$ for unit $\hat g_k$, so
$G\triangleq K I_3-\sum_k\hat g_k\hat g_k^{\!\top}\succeq 0$ and
\eqref{eq:closed} is a sum of two positive semidefinite terms. By
subadditivity of rank,
$\operatorname{rank}M\le\operatorname{rank}(vv^{\!\top})
+\operatorname{rank}(C^{\!\top}GC)\le 1+\operatorname{rank}C$, and
$C\in\mathbb{R}^{3\times P}$.
\end{IEEEproof}

\begin{corollary}[Rigid objects are capped at four]
\label{cor:rigid}
If the configuration is held fixed, then no number of gravity directions and no
number of repeated measurements can raise the rank above $4$. Consequently a
rigid assembly of $P\ge 5$ bodies is \emph{not} identifiable from quasi-static
wrench measurements without additional priors.
\end{corollary}

Corollary~\ref{cor:rigid} delimits what measurement-grounded methods that treat
the object as a rigid body can achieve, and explains why such methods must
introduce priors --- homogeneous-part assumptions, non-negativity, or
language-derived material guesses --- to disambiguate the remaining directions.
Our setting replaces those priors with measurements: changing $\theta$ changes
$C(\theta)$ and thereby injects genuinely new rows.

\begin{proposition}[What remains invisible]
\label{prop:null}
With an orthonormal triad,
%
\begin{equation}
\mathcal{N}\big(M(\theta)\big)=\Big\{u\in\mathbb{R}^P:\;
\textstyle\sum_i u_iV_i=0\;\;\text{and}\;\;\sum_i u_iV_ic_i(\theta)=0\Big\}.
\end{equation}
%
That is, a redistribution of mass that preserves both the total mass and the
first moment about the sensor origin is unobservable at $\theta$.
\end{proposition}

\begin{IEEEproof}
Write \eqref{eq:triad} as $M=a\,vv^{\!\top}+b\,C^{\!\top}C$ with $a,b>0$. For
$u\in\mathbb{R}^P$, $Mu=0\iff u^{\!\top}\!Mu=0
\iff a(v^{\!\top}u)^2+b\|Cu\|^2=0\iff v^{\!\top}u=0$ and $C(\theta)u=0$.
\end{IEEEproof}

Proposition~\ref{prop:null} is the exact statement of what active exploration
must undo: a redistribution hidden at $\theta_1$ becomes visible at $\theta_2$
precisely when $C(\theta_2)u\neq0$. Articulation is the only mechanism that
moves $C$, which is why the method is specific to articulated objects rather
than merely applicable to them.

\subsection{How Many Configurations Are Necessary}
\label{sec:rounds}

\begin{proposition}[Lower bound on rounds]
\label{prop:rounds}
After measuring at configurations $\theta_1,\dots,\theta_R$ with
$M_R=\sum_r M(\theta_r)$,
%
\begin{equation}
\operatorname{rank}M_R\;\le\;1+\min(3R,\,P),
\end{equation}
%
so identifying all $P$ densities requires
$R\ge\lceil (P-1)/3\rceil$ distinct configurations.
\end{proposition}

\begin{IEEEproof}
$\sum_r C(\theta_r)^{\!\top}C(\theta_r)=\mathcal{C}^{\!\top}\mathcal{C}$ where
$\mathcal{C}=[C(\theta_1)^{\!\top}\cdots C(\theta_R)^{\!\top}]^{\!\top}
\in\mathbb{R}^{3R\times P}$, whose rank is at most $\min(3R,P)$. Apply rank
subadditivity as in Proposition~\ref{prop:rank}.
\end{IEEEproof}

For a serial chain of $p$ links carrying a hinge at each of its $p-1$ joints,
the estimation vector has $P=2p-1$ entries and Proposition~\ref{prop:rounds}
gives
%
\begin{equation}
R\;\ge\;\Big\lceil \tfrac{2p-2}{3}\Big\rceil .
\label{eq:chain-bound}
\end{equation}
%
This reproduces the observed behaviour of our two custom objects exactly: the
two-link object ($p=2$, $P=3$) converges in one round and the three-link object
($p=3$, $P=5$) in two, matching the bound in both cases. Note that
\eqref{eq:chain-bound} is computed from geometry alone, since $M$ does not
depend on $\rho$ (Remark~\ref{rem:closed-loop}).

\begin{remark}[Necessary, not sufficient]
\label{rem:conditioning}
Proposition~\ref{prop:rounds} constrains rank only. Once the information matrix
is full rank the number of rounds actually required is governed by its
conditioning, which degrades sharply with the angle-measurement error because
that error perturbs $A$ itself rather than $y$ --- an errors-in-variables
structure under which ordinary least squares is asymptotically biased. We
therefore estimate $\rho$ by a total-least-squares formulation, and quantify
the gap between the rank bound and the achievable round count as a function of
angular precision in Section~\ref{sec:exp}.
\end{remark}
